\chapter{Tree algorithms: subtrees, diameters, LCA}\label{ch:trees}
\index{tree}\index{LCA}\index{binary lifting}\index{rerooting}

A tree is the graph where every question has a canonical answer:
exactly one path between any two vertices, and rooting turns
``around'' into ``down.'' This chapter collects the five patterns
that solve almost every tree problem on the archive: post-order
subtree sums, the two-DFS diameter, rerooting (answers for every
root in one pass), the Euler tour that turns subtrees into subarrays,
and binary lifting for ancestors and LCA. After it you will compute
per-vertex farthest distances, answer distance queries in
$O(\log n)$, and see every ``tree problem'' as pick-a-pattern.

% from-zero:begin
\begin{rulebox}[Before you start]
Chapter~7 (DFS) and Chapter~9 (a tree is a connected graph with no cycle). \emph{Recap:} in a tree there is exactly one simple path between any two vertices, which is why tree problems are easier than general graph problems.
\end{rulebox}
\begin{wins}[After this chapter you can]
\begin{itemize}
\item root a tree and compute subtree sizes, depths and sums with one DFS;
\item find a tree's diameter with two DFS passes;
\item flatten a tree into an array with an Euler tour so range tools apply;
\item answer lowest-common-ancestor and path-distance queries with binary lifting.
\end{itemize}
\end{wins}
\begin{principle}[How to think about this]
A tree is a graph with no loops, so there is exactly one way to get anywhere. That removes most of the difficulty: you can solve a problem for a small subtree and combine the answers upward. The habit to build is ``solve my children first, then combine,'' which is a recursive way of thinking that returns again in dynamic programming.
\end{principle}
% from-zero:end

\section{The shape of a tree}\label{sec:tshape}
$n$ vertices, $n{-}1$ edges, connected, acyclic: any two properties
imply the third. Rooting is a \emph{design choice} --- most
algorithms pick an arbitrary root and exploit ``parent vs children''
as the graph's only direction. Unique paths are what make ancestors,
depths and diameters well-defined.

\section{Subtree aggregates}\label{sec:subtree}\index{subtree aggregate}
Post-order accumulation computes any quantity that depends on
children first: sizes, sums, counts. The pattern is one DFS;
everything else in the chapter builds on it.
\begin{example}[CSES \emph{Subordinates}]\label{ex:subor}
Per employee, count the whole subtree beneath them.
\par\smallskip
\textbf{Strategy.} The hierarchy is already rooted: size(v) =
$\sum_{u\text{ child}}(1+\text{size}(u))$ computed in one post-order
pass.
\par\smallskip
\begin{solution}
\lstinputlisting{code/subordinates.cpp}
Real output (CSES sample):
\begin{lstlisting}[numbers=none]
4 0 2 0 0
\end{lstlisting}
\end{solution}
\end{example}

\section{Diameter: two DFS}\label{sec:diam}\index{diameter}
The farthest vertex from \emph{any} start is a diameter endpoint;
a second farthest-search from it finds the other end. Proof sketch:
let $u,v$ be a diameter and $a$ the farthest vertex from start $s$;
if $a$ were not an endpoint, a path past $u$ or $v$ would beat
$a$'s distance --- contradiction. (The common wrong proof --- ``it
works because it works on paths'' --- is flagged in the mistakes.)
\begin{example}[CSES \emph{Tree Diameter}]\label{ex:diam}
Longest path in the tree.
\par\smallskip
\textbf{Strategy.} Two searches, no DP: farthest-from-1, then
farthest-from-that.
\par\smallskip
\begin{solution}
\lstinputlisting{code/tree_diameter.cpp}
Real output (CSES sample):
\begin{lstlisting}[numbers=none]
3
\end{lstlisting}
\end{solution}
\end{example}
\begin{figure}[ht]\centering
% fig47-twodfs.tex -- Figure 10.x: two-DFS diameter. Path u1-u2-u3-u4-u5 plus leaf a on u2.
% fstyle.tex -- shared TikZ styles for the figures of this book (palette from lifeos.sty).
\tikzset{%
  cell/.style={draw=ink, minimum width=7mm, minimum height=6.5mm, inner sep=0pt, font=\footnotesize, fill=white},
  cellh/.style={cell, fill=accentsoft, draw=accent},
  cella/.style={cell, fill=amberbg, draw=amber},
  cellg/.style={cell, fill=goodbg, draw=good},
  cellx/.style={cell, fill=graybg, draw=slate, text=slate},
  vtx/.style={circle, draw=ink, minimum size=6.5mm, inner sep=0pt, font=\footnotesize, fill=white},
  vtxh/.style={vtx, draw=accent, fill=accentsoft},
  vtxa/.style={vtx, draw=amber, fill=amberbg},
  vtxg/.style={vtx, draw=good, fill=goodbg},
  vtxd/.style={vtx, draw=slate, fill=graybg, text=slate},
  lab/.style={font=\footnotesize, text=slate},
  labA/.style={font=\footnotesize, text=accent},
  labM/.style={font=\footnotesize, text=amber},
  idx/.style={font=\tiny, text=slate},
  arr/.style={-{stealth}, thick, ink},
  arrA/.style={-{stealth}, thick, accent},
  arrM/.style={-{stealth}, thick, amber},
  arrG/.style={-{stealth}, thick, good},
  edge/.style={thick, ink},
  edgeA/.style={thick, accent},
  edgeM/.style={thick, amber},
}%
% \cellrow{prefix}{x0}{y}{style}{v0,v1,...}: a row of cells named prefix0, prefix1, ... (pitch 0.75cm)
\providecommand{\cellrow}[5]{\foreach \v [count=\i from 0] in {#5} {\node[#4] (#1\i) at (#2+\i*0.75,#3) {\v};}}%
% \cellidx{prefix}{n}{start}: index labels under cells prefix0 .. prefix(n-1)
\providecommand{\cellidx}[3]{\foreach \i in {0,...,\numexpr#2-1\relax} {\node[idx, below=0pt of #1\i] {\the\numexpr\i+#3\relax};}}%
% \fitfig{<tikzpicture>}: shrink the picture to \linewidth only when it would overflow.
\ifdefined\fitbox\else\newsavebox{\fitbox}\fi%
\providecommand{\fitfig}[1]{\sbox{\fitbox}{#1}\ifdim\wd\fitbox>\linewidth\resizebox{\linewidth}{!}{\usebox{\fitbox}}\else\usebox{\fitbox}\fi}%

\begin{tikzpicture}[font=\footnotesize]
\foreach \i in {1,...,5} \node[vtx] (u\i) at (\i*1.6,0) {$u_\i$};
\node[vtxh] (a) at (3.2,1.4) {$a$};
\draw[edge] (u1)--(u2)--(u3)--(u4)--(u5) (a)--(u2);
\node[vtxa, minimum size=6.5mm] at (u5) {$u_5$};
\draw[edgeM, very thick] (a)--(u2)--(u3)--(u4)--(u5);
\draw[edgeA, very thick, dashed, transform canvas={yshift=-3pt}] (u5)--(u4)--(u3)--(u2)--(u1);
\node[labM, anchor=west] at (0.7,-1.0) {1st DFS from $a$ (amber path): farthest vertex is $u_5$, distance $4$};
\node[labA, anchor=west] at (0.7,-1.6) {2nd DFS from $u_5$ (blue dashes): farthest is $u_1$, distance $4$ $=$ the diameter};
\end{tikzpicture}

\caption{The farthest-point argument. Notice the claim is about
endpoints, not the path --- endpoint first, distance second.}
\label{fig:twodfs}
\end{figure}

\section{Rerooting}\label{sec:reroot}\index{rerooting}
``For every root'' doubles the work naively; rerooting does it in
two passes: \texttt{down[v]} from children, then \texttt{up[v]}
passed down from parent using each child's best \emph{sibling}
value. Per vertex the answer combines both.
\begin{example}[CSES \emph{Tree Distances II}]\label{ex:dist2}
Each vertex's farthest distance.
\par\smallskip
\textbf{Strategy.} The answer is $\max(\texttt{down},\texttt{up})$;
the sibling-max trick supplies each child's ``up'' in $O(\deg v)$.
\par\smallskip
\begin{solution}
\lstinputlisting{code/tree_distances2.cpp}
Real output (CSES sample):
\begin{lstlisting}[numbers=none]
2 3 2 3 3
\end{lstlisting}
\end{solution}
\end{example}
\begin{figure}[ht]\centering
% fig48-reroot.tex -- Figure 10.x: rerooting; up[u] = 1 + max(up[v], 1 + down[w]) for sibling w. Here up[v]=1, down[w]=2.
% fstyle.tex -- shared TikZ styles for the figures of this book (palette from lifeos.sty).
\tikzset{%
  cell/.style={draw=ink, minimum width=7mm, minimum height=6.5mm, inner sep=0pt, font=\footnotesize, fill=white},
  cellh/.style={cell, fill=accentsoft, draw=accent},
  cella/.style={cell, fill=amberbg, draw=amber},
  cellg/.style={cell, fill=goodbg, draw=good},
  cellx/.style={cell, fill=graybg, draw=slate, text=slate},
  vtx/.style={circle, draw=ink, minimum size=6.5mm, inner sep=0pt, font=\footnotesize, fill=white},
  vtxh/.style={vtx, draw=accent, fill=accentsoft},
  vtxa/.style={vtx, draw=amber, fill=amberbg},
  vtxg/.style={vtx, draw=good, fill=goodbg},
  vtxd/.style={vtx, draw=slate, fill=graybg, text=slate},
  lab/.style={font=\footnotesize, text=slate},
  labA/.style={font=\footnotesize, text=accent},
  labM/.style={font=\footnotesize, text=amber},
  idx/.style={font=\tiny, text=slate},
  arr/.style={-{stealth}, thick, ink},
  arrA/.style={-{stealth}, thick, accent},
  arrM/.style={-{stealth}, thick, amber},
  arrG/.style={-{stealth}, thick, good},
  edge/.style={thick, ink},
  edgeA/.style={thick, accent},
  edgeM/.style={thick, amber},
}%
% \cellrow{prefix}{x0}{y}{style}{v0,v1,...}: a row of cells named prefix0, prefix1, ... (pitch 0.75cm)
\providecommand{\cellrow}[5]{\foreach \v [count=\i from 0] in {#5} {\node[#4] (#1\i) at (#2+\i*0.75,#3) {\v};}}%
% \cellidx{prefix}{n}{start}: index labels under cells prefix0 .. prefix(n-1)
\providecommand{\cellidx}[3]{\foreach \i in {0,...,\numexpr#2-1\relax} {\node[idx, below=0pt of #1\i] {\the\numexpr\i+#3\relax};}}%
% \fitfig{<tikzpicture>}: shrink the picture to \linewidth only when it would overflow.
\ifdefined\fitbox\else\newsavebox{\fitbox}\fi%
\providecommand{\fitfig}[1]{\sbox{\fitbox}{#1}\ifdim\wd\fitbox>\linewidth\resizebox{\linewidth}{!}{\usebox{\fitbox}}\else\usebox{\fitbox}\fi}%

\begin{tikzpicture}[font=\footnotesize]
\node[vtx] (p) at (2,3.6) {$p$}; \node[vtx] (v) at (2,2.4) {$v$};
\node[vtxh] (u) at (0.8,1.2) {$u$}; \node[vtx] (w) at (3.2,1.2) {$w$};
\node[vtx] (x) at (3.2,0) {$x$}; \node[vtx] (y) at (3.2,-1.2) {$y$};
\draw[edge] (p)--(v)--(u) (v)--(w)--(x)--(y);
\draw[arrA, bend right=40] (p.west) to node[left, labA] {up$[v]=1$} (v.west);
\draw[arrM, bend left=35, dashed] (w.east) to node[right, labM, align=left] {down$[w]=2$} (y.east);
\node[labM, align=left, anchor=west] at (6.3,1.9) {pass 2: up$[u]$ $=1+\max\big(\mathrm{up}[v],\ 1+\mathrm{down}[w]\big)$};
\node[labM, align=left, anchor=west] at (6.3,1.2) {\phantom{pass 2: up$[u]$ }$=1+\max(1,3)=4$};
\node[lab, align=left, anchor=west, text width=6.2cm] at (6.3,0.2) {the longest path from $u$ goes up to $v$, then down through the \emph{sibling} $w$ to $y$};
\end{tikzpicture}

\caption{What travels down in pass two. Notice a child never sees its
own subtree in its ``up'' value --- the sibling max is the key step.}
\label{fig:reroot}
\end{figure}

\section{The Euler tour}\label{sec:euler}\index{Euler tour}
\texttt{tin[v]} on entry, \texttt{tout[v]} on exit: every subtree
becomes the interval $[\texttt{tin},\texttt{tout})$, turning
``update/query a subtree'' into range operations --- the bridge to
segment trees in Chapter~12.
\begin{figure}[ht]\centering
\begin{tikzpicture}[font=\footnotesize]
\node[circle,draw=slate,font=\tiny] (r) at (2,1.6) {1};
\node[circle,draw=slate,font=\tiny] (a) at (1,0.8) {2};
\node[circle,draw=slate,font=\tiny] (b) at (3,0.8) {3};
\node[circle,draw=slate,font=\tiny] (c) at (1,0) {4};
\draw[slate] (r)--(a) (r)--(b) (a)--(c);
\node[font=\tiny,anchor=west] at (4.2,1.2) {tin: 1=0, 2=1, 4=2, 3=3};
\node[font=\tiny,anchor=west] at (4.2,0.7) {tout: 1=4, 2=3, 4=3, 3=4};
\node[font=\tiny,anchor=west] at (4.2,0.2) {subtree(2) = $[1,3)$: $\{2,4\}$};
\node[font=\tiny] at (2,-0.8) {every subtree is one interval};
\end{tikzpicture}

\caption{Euler intervals nested. Notice containment is interval
containment --- $u$ is in $v$'s subtree iff
$\texttt{tin}[v]\le\texttt{tin}[u]<\texttt{tout}[v]$.}
\label{fig:euler}
\end{figure}

\section{Binary lifting and LCA}\label{sec:lca}\index{binary lifting}
Store the $2^k$-th ancestor of every vertex: $\texttt{up}[k][v] =
\texttt{up}[k{-}1][\texttt{up}[k{-}1][v]]$ --- a doubling table built
in $O(n\log n)$. $k$-th ancestor: follow set bits. LCA: equalize
depth, then lift both ends by descending powers of two until their
parents coincide. (An Euler-tour + sparse-RMQ alternative exists;
lifting is the workhorse.)
\begin{figure}[ht]\centering
\begin{tikzpicture}[font=\footnotesize]
\foreach \i/\v in {0/7,1/5,2/3,3/1}
  \node[circle,draw=slate,font=\tiny] (\v) at (\i*1.1,0) {\v};
\draw[slate] (7)--(5) (5)--(3) (3)--(1);
\draw[-{stealth},amber,thick] (7) to[out=30,in=150] node[above,font=\tiny] {$2^1$} (3);
\draw[-{stealth},amber,thick] (7) to[out=45,in=135] node[above,font=\tiny] {$2^2$} (1);
\node[font=\tiny] at (1.7,-0.9) {$up[v][k]$ jumps $2^k$ ancestors};
\end{tikzpicture}

\caption{The doubling table. Notice each row skips exactly $2^k$
ancestors --- a jump of $2^k$ is two jumps of $2^{k-1}$.}
\label{fig:binlift}
\end{figure}
\begin{lstlisting}[language=C++,caption={\texttt{lca.h}: lifting
table, \texttt{kth}, \texttt{lca}, \texttt{dist}.},label={lst:lca}]
\end{lstlisting}
\lstinputlisting{code/lca.h}
\begin{example}[CSES \emph{Company Queries I \& II}]\label{ex:comp}
$k$-th boss; lowest common boss.
\par\smallskip
\textbf{Strategy.} The hierarchy is rooted at the CEO: \texttt{kth}
answers I, \texttt{lca} answers II --- same table.
\par\smallskip
\begin{solution}
\lstinputlisting{code/company_queries.cpp}
Real output (CSES sample):
\begin{lstlisting}[numbers=none]
3
1
-1
\end{lstlisting}
\end{solution}
\end{example}

\section{Distance queries}\label{sec:distq}
$\texttt{dist}(u,v) = \texttt{depth}[u] + \texttt{depth}[v] -
2\cdot\texttt{depth}[\texttt{lca}(u,v)]$ --- the unique path is the
only route, so depth arithmetic is exact.
\begin{example}[CSES \emph{Distance Queries}]\label{ex:distq}
Pairwise distances on a tree.
\par\smallskip
\textbf{Strategy.} LCA plus depth, $O(\log n)$ per query.
\par\smallskip
\begin{solution}
\lstinputlisting{code/distance_queries.cpp}
Real output (CSES sample):
\begin{lstlisting}[numbers=none]
2
3
0
\end{lstlisting}
\end{solution}
\end{example}
\begin{figure}[ht]\centering
\begin{tikzpicture}[font=\footnotesize]
\node[font=\tiny,anchor=west] at (0,0.6) {lift deeper end to equal depth};
\node[font=\tiny,anchor=west] at (0,0.1) {binary-search both lifts together};
\draw[-{stealth},slate] (0,-0.4)--(0,-1.1);
\node[font=\tiny,anchor=west] at (0,-1.5) {parents meet $\Rightarrow$ LCA};
\end{tikzpicture}

\caption{LCA in two lifts. Notice the first lift is exact depth
equalization; the second is the binary search that stops one below
the meeting point.}
\label{fig:lca}
\end{figure}

\begin{exercise}[easy]\label{exr:rooted}
When does an unrooted answer (diameter) differ from a rooted one
(deepest child)? Name a quantity insensitive to rooting.
\end{exercise}
\begin{exercise}[medium]\label{exr:farthest}
Complete the farthest-vertex proof: if $a$'s path to a diameter
endpoint crosses the diameter at $x$, which distance comparison
contradicts $a$ not being an endpoint?
\end{exercise}
\begin{exercise}[medium]\label{exr:liftfill}
Fill the lift-table recurrence
$\texttt{up}[k][v]=\texttt{up}[k{-}1][\texttt{up}[k{-}1][v]]$ and
argue it counts exactly $2^k$ edges.
\end{exercise}
\begin{exercise}[medium]\label{exr:subarray}
Convert ``add 1 to every vertex in $v$'s subtree'' into a
subarray operation using \texttt{tin}/\texttt{tout}.
\end{exercise}

\section{Chapter summary}
\begin{table}[ht]\centering
{\small
\begin{tabularx}{\linewidth}{@{}lYYY@{}}
\toprule
\tablehead\thead{Need} & \thead{Tool} & \thead{Cost} & \thead{Yields}\\
\midrule
subtree sums & post-order & $O(n)$ & aggregates\\
longest path & two-DFS & $O(n)$ & diameter\\
answer per root & reroot & $O(n)$ & down+up\\
$k$-th anc, LCA & lifting & $O(n\log n)$+$O(\log n)$/q & dist\\
subtree as range & Euler & $O(n)$ & range ops\\
\bottomrule
\end{tabularx}}
\caption{Root it, then pick the pattern: aggregates below, lifts
above, intervals across.}
\label{tab:trees}
\end{table}

\section{Common mistakes}
\index{common mistakes}
\begin{itemize}
\item \textbf{The lazy diameter ``proof.}'' ``Works on paths'' is not
a proof --- the farthest-point argument is required and is exactly
what breaks if you skip it.
\item \textbf{Recursion depth on a path graph.} $n=2\cdot10^5$ chain
overflows a default stack --- convert to iterative or
\texttt{vector}-stack DFS.
\item \textbf{Lifting with \texttt{up[k][v]} unguarded.} Root's
ancestors are 0; dereferencing \texttt{up[.][0]} must be defined
(index-0 sentinel) not skipped ad hoc.
\item \textbf{Sibling-max off-by-one.} The child's own \texttt{down}
must be \emph{excluded} from its up-value --- prefix/suffix maxima
around its index, not the global max.
\item \textbf{Euler half-open confusion.} \texttt{tout} is exclusive
on exit-time --- mixing $[$tin,tout$]$ and $[$tin,tout$)$ silently
includes the next vertex.
\end{itemize}

\section{Additional exercises}
\begin{exercise}[easy] CSES \emph{Tree Diameter}. \end{exercise}
\begin{exercise}[easy] CSES \emph{Tree Distances I} --- BFS both ends. \end{exercise}
\begin{exercise}[easy] CSES \emph{Subordinates}. \end{exercise}
\begin{exercise}[easy] USACO-S \emph{MooTube} --- threshold component sizes. \end{exercise}
\begin{exercise}[easy] USACO-B \emph{Family Tree} --- LCA-style ancestry. \end{exercise}
\begin{exercise}[medium] CSES \emph{Tree Matching} --- greedy DP on children. \end{exercise}
\begin{exercise}[medium] CSES \emph{Company Queries II} --- LCA direct. \end{exercise}
\begin{exercise}[medium] USACO-S \emph{Milk Visits} --- LCA + path membership. \end{exercise}
\begin{exercise}[medium] AC \emph{Count Descendants} --- Euler + count. \end{exercise}
\begin{exercise}[medium] CF \emph{Minimize the Diameter} --- reroot variant. \end{exercise}
\begin{exercise}[medium] CSES \emph{Subtree Queries} --- Euler tour + range structure. \end{exercise}
\begin{exercise}[medium] CSES \emph{Path Queries} --- the same tour answering root-to-node aggregates. \end{exercise}
\begin{exercise}[hard] CSES \emph{Finding a Centroid} --- subtree-size bisect. \end{exercise}
\begin{exercise}[hard] CSES \emph{Counting Paths} --- difference array on paths. \end{exercise}
\begin{exercise}[hard] CSES \emph{Fixed-Length Paths I} --- centroid decomposition preview. \end{exercise}

\chapternotes
Post-order/reroot/diameter patterns follow Laaksonen's tree
chapter~\cite{laaksonen2024} and Sannemo~\cite{sannemo2025}; binary
lifting and Euler-tour RMQ are the USACO Guide's Gold tree
modules~\cite{usacoguide}; cp-algorithms carries the two-DFS proof
and lifting table~\cite{cpalgorithms}; Halim supplies named
variants~\cite{halimcp4}; archive problems are official
CSES/USACO~\cite{cses}.

\section*{Further reading}
\begin{itemize}
\item USACO Guide, ``Introduction to Tree Algorithms''
(intermediate) ---
subtree aggregations and the two-DFS diameter of \S\ref{sec:diam}. \newline\url{https://usaco.guide/silver/intro-tree}
\item USACO Guide, ``Euler Tour Technique'' (advanced) --- the tour view of
\S\ref{sec:euler} with data-structure queries on top. \newline\url{https://usaco.guide/gold/tree-euler}
\item USACO Guide, ``Binary Jumping'' (advanced) --- the $2^k$ ancestors
machinery behind \S\ref{sec:lca}/\S\ref{sec:distq}. \newline\url{https://usaco.guide/plat/binary-jump}
\item cp-algorithms, ``Lowest Common Ancestor --- Binary Lifting''
(reference) ---
\S\ref{sec:lca} with the $\mathrm{up}[k][v]$ table in full. \newline\url{https://cp-algorithms.com/graph/lca_binary_lifting.html}
\item CSES 1131 \emph{Tree Diameter} and 1135 \emph{Distance Queries}
(practice) ---
\ref{ex:diam}/\ref{ex:distq} verbatim. \newline\url{https://cses.fi/problemset/task/1131}
\item CSES 1688 \emph{Company Queries II} (practice) --- LCA-as-submission for
\ref{ex:comp}. \newline\url{https://cses.fi/problemset/task/1688}\end{itemize}

\begin{rulebox}[Watch instead]
\smallskip\noindent\textbf{Best first.}
``LCA -- Lowest Common Ancestor'' --- Errichto Algorithms, 16 min ---
\url{https://youtu.be/dOAxrhAUIhA}. The binary-lifting
walk of \S\ref{sec:lca}--\S\ref{sec:distq} from a contestant's angle;
same \texttt{up[k][v]} table, slightly different lift order in code.
\smallskip\noindent
``How to find Diameter of a Tree'' --- Luv, 12 min ---
\url{https://youtu.be/m8qXEz1Xe8w}. The two-DFS
diameter trick of \S\ref{sec:diam} plus why it is correct on trees
only --- matches the book's caution about weighted/non-tree inputs.
Access date 2026-10-10.
\end{rulebox}
